\section{Quasi-compact morphisms}
\label{section-quasi-compact}

\noindent
A scheme is {\it quasi-compact} if its underlying topological space is
quasi-compact. There is a relative notion which is defined as follows.

\begin{definition}
\label{definition-quasi-compact}
A morphism of schemes is called {\it quasi-compact}
if the underlying map of topological spaces is
quasi-compact, see
Topology, Definition \ref{topology-definition-quasi-compact}.
\end{definition}

\begin{lemma}
\label{lemma-quasi-compact-affine}
Let $f : X \to S$ be a morphism of schemes.
The following are equivalent
\begin{enumerate}
\item $f : X \to S$ is quasi-compact,
\item the inverse image of every affine open is quasi-compact, and
\item there exists some affine open covering $S = \bigcup_{i \in I} U_i$
such that $f^{-1}(U_i)$ is quasi-compact for all $i$.
\end{enumerate}
\end{lemma}

\begin{proof}
Suppose we are given a covering $S = \bigcup_{i \in I} U_i$ as in (3).
First, let $U \subset S$ be any affine open. For any $u \in U$
we can find an index $i(u) \in I$ such that $u \in U_{i(u)}$.
As standard opens form a basis for the topology on $U_{i(u)}$ we can find
$W_u \subset U \cap U_{i(u)}$ which is standard open in $U_{i(u)}$.
By compactness we can find finitely many points $u_1, \ldots, u_n \in U$
such that $U = \bigcup_{j = 1}^n W_{u_j}$. For each $j$ write
$f^{-1}U_{i(u_j)} = \bigcup_{k \in K_j} V_{jk}$ as a finite
union of affine opens. Since $W_{u_j} \subset U_{i(u_j)}$ is a standard
open we see that $f^{-1}(W_{u_j}) \cap V_{jk}$ is a standard
open of $V_{jk}$, see Algebra, Lemma \ref{algebra-lemma-spec-functorial}.
Hence $f^{-1}(W_{u_j}) \cap V_{jk}$ is affine, and so
$f^{-1}(W_{u_j})$ is a finite union of affines. This proves that the
inverse image of any affine open is a finite union of affine opens.

\medskip\noindent
Next, assume that the inverse image of every affine open is a finite
union of affine opens.
Let $K \subset S$ be any quasi-compact open. Since $S$ has a basis
of the topology consisting of affine opens we see that $K$ is a finite
union of affine opens. Hence the inverse image of $K$ is a finite
union of affine opens. Hence $f$ is quasi-compact.

\medskip\noindent
Finally, assume that $f$ is quasi-compact. In this case the argument
of the previous paragraph shows that the inverse image of any affine
is a finite union of affine opens.
\end{proof}

\begin{lemma}
\label{lemma-quasi-compact-preserved-base-change}
Being quasi-compact is a property of morphisms of schemes
over a base which is preserved under arbitrary base change.
\end{lemma}

\begin{proof}
Omitted.
\end{proof}

\begin{lemma}
\label{lemma-composition-quasi-compact}
The composition of quasi-compact morphisms is quasi-compact.
\end{lemma}

\begin{proof}
This follows from the definitions and
Topology, Lemma \ref{topology-lemma-composition-quasi-compact}.
\end{proof}

\begin{lemma}
\label{lemma-closed-immersion-quasi-compact}
A closed immersion is quasi-compact.
\end{lemma}

\begin{proof}
Follows from the definitions and
Topology, Lemma \ref{topology-lemma-closed-in-quasi-compact}.
\end{proof}

\begin{example}
\label{example-open-immersion-not-quasi-compact}
An open immersion is in general not quasi-compact.
The standard example of this is the open subspace
$U \subset X$, where $X = \Spec(k[x_1, x_2, x_3, \ldots])$,
where $U$ is $X \setminus \{0\}$, and where $0$ is the point
of $X$ corresponding to the maximal ideal
$(x_1, x_2, x_3, \ldots)$.
\end{example}

\begin{lemma}
\label{lemma-image-quasi-compact-closed}
Let $f : X \to S$ be a quasi-compact morphism of schemes.
The following are equivalent
\begin{enumerate}
\item $f(X) \subset S$ is closed, and
\item $f(X) \subset S$ is stable under specialization.
\end{enumerate}
\end{lemma}

\begin{proof}
We have (1) $\Rightarrow$ (2) by
Topology, Lemma \ref{topology-lemma-open-closed-specialization}.
Assume (2). Let $U \subset S$ be an affine open. It suffices to prove
that $f(X) \cap U$ is closed. Since $U \cap f(X)$ is stable under
specializations in $U$, we have reduced to the case where $S$ is affine.
Because $f$ is quasi-compact we deduce that $X = f^{-1}(S)$ is
quasi-compact as $S$ is affine. Thus we may write
$X = \bigcup_{i = 1}^n U_i$ with $U_i \subset X$ open affine.
Say $S = \Spec(R)$ and
$U_i = \Spec(A_i)$ for some $R$-algebra $A_i$.
Then $f(X) = \Im(\Spec(A_1 \times \ldots \times A_n)
\to \Spec(R))$. Thus the lemma follows from
Algebra, Lemma \ref{algebra-lemma-image-stable-specialization-closed}.
\end{proof}

\begin{lemma}
\label{lemma-quasi-compact-closed}
Let $f : X \to S$ be a quasi-compact morphism of schemes.
Then $f$ is closed if and only if specializations lift
along $f$, see
Topology, Definition \ref{topology-definition-lift-specializations}.
\end{lemma}

\begin{proof}
According to
Topology, Lemma \ref{topology-lemma-closed-open-map-specialization}
if $f$ is closed then specializations lift along $f$.
Conversely, suppose that specializations lift along $f$.
Let $Z \subset X$ be a closed subset. We may think of $Z$
as a scheme with the reduced induced scheme structure, see
Definition \ref{definition-reduced-induced-scheme}.
Since $Z \subset X$ is closed the restriction
of $f$ to $Z$ is still quasi-compact. Moreover specializations lift
along $Z \to S$ as well,
see Topology, Lemma \ref{topology-lemma-lift-specialization-composition}.
Hence it suffices to prove $f(X)$ is closed if specializations lift along $f$.
In particular $f(X)$ is stable under specializations, see
Topology, Lemma \ref{topology-lemma-lift-specializations-images}.
Thus $f(X)$ is closed by
Lemma \ref{lemma-image-quasi-compact-closed}.
\end{proof}










